% !PRINTED 281
\weberpage{281}
%<<CONT
a plane curve of the third order, if one sets $x^8=X$,
$y^8=Y$:

\weq{3}{X\,Y(X-Y)=16.}

If we set [§ 75, (1)]

\weq{4}{\xi=f\left(\dfrac{c+\partial\omega}{a}\right),\ \eta=\left(\dfrac{2}{a}\right)f_1\left(\dfrac{c+\partial\omega}{a}\right),\ \zeta=\left(\dfrac{2}{\partial}\right)f_2\left(\dfrac{c+\partial\omega}{a}\right),}

then the quantities $\xi, \eta, \zeta$ likewise satisfy the equations (2), since % [sic] German source prints the singular noun and verb for three quantities.
\[
\left(\dfrac{2}{a}\right)\left(\dfrac{2}{\partial}\right)=\left(\dfrac{2}{n}\right)=+1
\]
holds, and the point $\Pi$, whose coordinates are $\xi, \eta, \zeta$, therefore also lies
on the fundamental curve.

An equation
\[
\Phi\,(x,\,y,\,z,\,\xi,\,\eta,\,\zeta)=0
\]
% [sic] The German source calls this equation (5), but prints no number here and assigns (5) to the following display.
denotes, when the point $P$ is held fixed, a \gesperrt{surface} on
which $\Pi$ is to lie, and the intersection of this surface with the
fundamental curve gives a certain number of points $\Pi$ that correspond to the
point $P$. In the same way, to a fixed point $\Pi$ there corresponds
a certain number of points $P$, and this assignment is called
a \gesperrt{correspondence on the fundamental curve.} If the point $P$
coincides with one of the points $\Pi$ corresponding to it, we obtain a \gesperrt{double point or coincidence of the correspondence.}

If equation (5) is not altered by the interchange of $P$ with $\Pi$,
the correspondence is a \gesperrt{mutual one}.
Such correspondences are given by equations (8), (9),
(10), § 75.

By § 75, (3) we have then to set

\weq{5}{\begin{aligned} 2A &= x\xi+(-1)^{\frac{n+1}{8}}(y\eta+z\zeta)\\ B &= x\xi y\eta+x\xi z\zeta+(-1)^{\frac{n+1}{8}}y\eta z\zeta \end{aligned}}

$A$ is of the first, $B$ of the second degree with respect to
the coordinates of each of the two points $P$ and $\Pi$.

We wish to determine the degree of these correspondences, that is, the
number of points $\Pi$ corresponding to a point $P$. In doing so we have to multiply the degree of the equation $\Phi=0$ with
respect to $\xi, \eta, \zeta$ by the degree of the fundamental curve, that is, by 24,
and there results:
%CONT>>
